% =====================================================================
\chapter{Les sources de données et leur intégration}
\label{ch:donnees}
% =====================================================================

\typhoon{} n'invente aucune donnée : il assemble des sources publiques. La
qualité du produit dépend donc entièrement de la qualité de cette intégration.
Ce chapitre décrit chaque source, ce qu'elle apporte, ses contraintes
juridiques et techniques, et les pièges rencontrés.

\section{Vue d'ensemble}

\begin{table}[H]
\centering
\footnotesize
\caption{Les sources de données intégrées, leur rôle et leur niveau de
résolution.}
\label{tab:sources}
\begin{tabularx}{\textwidth}{@{}l l l X@{}}
\toprule
\textbf{Source} & \textbf{Rôle} & \textbf{Accès} & \textbf{Résolution} \\
\midrule
Géorisques (WFS) & Périmètres d'aléas & OGC WFS 2.0, GML 3.2 &
  \code{per-building} par test point-in-polygon \\
Géorisques (REST) & Statut communal, arrêtés CatNat & JSON &
  \code{commune-level} \\
BDNB (CSTB) & Vulnérabilité du bâtiment & REST JSON &
  \code{per-building} (identifiant bâtiment) \\
Géoplateforme (IGN) & Géocodage & REST JSON &
  Point + code INSEE \\
Classes TRI & Profondeur d'eau réglementaire & REST JSON &
  Par point, par classe de retour \\
Sandre & Tracé hydrographique & REST JSON & Tracé \\
Open-Meteo & Pluie horaire, débit GloFAS & REST JSON & Point, horaire \\
Panoramax & Photographie terrain & REST JSON & Point, orientée \\
Copernicus CDS & Projections climatiques 2050/2100 & API CDS & Grille $\approx$ 10~km \\
Mistral & Rédaction du rapport & REST & — (présentation) \\
\bottomrule
\end{tabularx}
\end{table}

\section{Géorisques : le cœur de la jointure}

\subsection{Ce que la source fournit}

Géorisques expose le référentiel des aléas réglementaires français. Il sert
deux interfaces distinctes, et cette dualité est structurante :

\begin{itemize}[leftmargin=1.4em]
  \item une interface \textbf{REST} par code INSEE, qui donne le statut
        communal : \og la commune est-elle concernée par un risque
        inondation ? \fg{} ;
  \item une interface \textbf{WFS} (Web Feature Service) qui expose les
        géométries vectorielles, donc permet de tester un point précis contre
        un polygone.
\end{itemize}

\subsection{Le test point-in-polygon}

Le WFS renvoie du GML 3.2, un dialecte XML volumineux. Le choix a été d'écrire
le parseur et le test géométrique à la main, sans dépendance à
\texttt{shapely}. Deux raisons : éviter une dépendance de compilation lourde
dans un contexte serverless, et garder la maîtrise complète des cas limites.

Le test utilise un \emph{lancer de rayon} (\emph{ray casting}) : on compte les
intersections d'un rayon partant du point avec les arêtes du polygone. Un
nombre impair signifie que le point est à l'intérieur. Le premier anneau d'un
polygone est l'anneau extérieur ; les suivants sont des trous, qu'il faut donc
soustraire.

\begin{lstlisting}[caption={Signature du test géométrique et gestion des
géométries non polygonales
(\fichier{backend/app/connectors/georisques\_wfs.py}).},label={lst:pip}]
def point_in_polygon(pt: tuple[float, float], polygon: Any) -> bool:
    """Point (lon, lat) dans un Polygon / MultiPolygon GML.

    L'anneau exterieur est le premier ; les suivants sont des trous.
    """

def geometry_intersects_point(
    geom: dict[str, Any],
    pt: tuple[float, float],
    radius_m: float = PROXIMITY_RADIUS_M,
) -> bool:
    """Le point (lon, lat) touche-t-il la geometrie ?

    Polygones -> point-in-polygon strict. Points/lignes -> test de proximite
    (une canalisation a 30 m est un risque pertinent meme si le batiment
    n'est pas dessus).
    """
\end{lstlisting}

\subsection{Les treize aléas}

Le référentiel est normalisé en treize objets aléa, chacun portant les mêmes
champs obligatoires :

\begin{center}
\small
\code{icpe}, \code{inondation}, \code{sismicite}, \code{mouvement\_terrain},
\code{radon}, \code{rga}, \code{cavite}, \code{feu\_foret}, \code{avalanche},
\code{canalisations}, \code{vent\_cyclonique}, \code{ppr}, \code{ssp}
\end{center}

\subsection{La ventilation par type de PPR}

Un piège est apparu ici. L'agrégat \code{ppr} répond à la question \og un plan
de prévention quelconque touche-t-il ce point ? \fg. C'est insuffisant pour un
travail par péril : un bâtiment peut être dans un périmètre PPR sismique sans
être dans un périmètre PPR inondation, et l'agrégat ne fait pas la différence.

Une table de correspondance a donc été ajoutée, qui ventile les huit couches
PPR par type de péril :

\begin{lstlisting}[caption={Ventilation des couches PPR par type de péril
(\fichier{georisques\_wfs.py}).},label={lst:ppr}]
# Ventilation des 8 couches PPR de WFS_LAYER_MAP["ppr"] par type de peril.
# L'agregat "ppr" ci-dessus repond "un PPR quelconque touche-t-il ce point ?"
# -- utile pour une carte generale, mais pas pour un peril precis.
PPR_TYPE_LAYERS: dict[str, list[str]] = {
    "inondation":        ["ms:PPRN_PERIMETRE_INOND", "ms:PPRN_PERIMETRE_SUBMAR"],
    "mouvement_terrain": ["ms:PPRN_PERIMETRE_MVT"],
    "seisme":            ["ms:PPRN_PERIMETRE_SEISME"],
    "avalanche":         ["ms:PPRN_PERIMETRE_AVALANCHE"],
    ...
}
\end{lstlisting}

\subsection{Le piège de la boîte englobante}

C'est le bug le plus instructif du projet. Pour interroger un WFS, il faut
fournir une \emph{boîte englobante} (BBOX). Le choix naturel est de prendre une
boîte \og raisonnable \fg, à l'échelle d'un quartier.

Erreur. Le WFS applique un plafond au nombre de géométries renvoyées
(paramètre \code{count}). Mesuré à Nantes : une boîte large se remplit de
polygones situés à 2~km, et le polygone contenant réellement l'adresse
\emph{n'arrive jamais dans la page de résultats}. Le test conclut alors \og
hors TRI \fg{} sur un quai inondable.

La correction a été de réduire la boîte à l'échelle du point : environ 55~m de
marge, soit \code{POINT\_\allowbreak MARGIN\_\allowbreak DEG}. Tout polygone
contenant le point coupe
nécessairement une boîte centrée sur lui : élargir n'ajoute aucune information
et fait perdre des réponses.

\begin{honesty}
Un résultat faux n'est pas toujours le signe d'un bug visible. Ici, le service
répondait correctement, sans erreur, avec une information fausse. C'est le
scénario le plus dangereux pour un produit de données : la panne silencieuse.
\end{honesty}

\subsection{Le piège du client HTTP partagé}

Autre découverte contre-intuitive : le WFS Géorisques exige un client
\code{httpx.AsyncClient} \emph{dédié}. La passerelle lie la connexion
persistante (\emph{keep-alive}) à un backend particulier. Mélanger les appels
REST et WFS sur le même client provoque des 404 en cascade — sans rapport
apparent avec la cause. Cette contrainte est documentée comme invariant et
vérifiée par un test d'intégration qui enchaîne REST puis WFS.

\subsection{Une réponse vide est ambiguë}

Une réponse WFS vide peut signifier deux choses : \og aucun périmètre ne
contient ce point \fg{} (fait) ou \og le serveur a hoqueté \fg{} (panne). La
distinction est essentielle pour le produit. Le connecteur rejoue donc la
requête une fois avant de conclure à une absence. C'est un compromis explicite :
un surcoût réseau en échange d'une réduction des faux négatifs.

\subsection{Les arrondissements}

Paris, Lyon et Marseille ont un découpage administratif particulier : les
arrondissements ont leur propre code INSEE (75101 à 75120 pour Paris), mais les
endpoints communaux attendent le code de la commune parente (75056). Sans
retraduction, la jointure échoue silencieusement. La correspondance est
implémentée et testée.

\section{La BDNB : la vulnérabilité}

La Base de Données Nationale des Bâtiments, produite par le CSTB, décrit chaque
bâtiment français. Le connecteur récupère environ 139 champs : usage, hauteur,
emprise au sol, matériaux de construction, année de construction, performance
énergétique, et des indicateurs de fiabilité de l'adressage.

Deux points de conception :

\begin{enumerate}[leftmargin=1.6em]
  \item \textbf{Passage verbatim.} Aucun renommage, aucun filtrage. Un test de
        conformité compare l'ensemble des clés reçues et l'ensemble des clés
        émises : ils doivent être identiques.
  \item \textbf{Ambiguïté d'adresse.} Plusieurs bâtiments peuvent partager une
        adresse. La première ligne renvoyée est traitée comme le bâtiment
        principal ; les autres sont conservées verbatim dans
        \code{autres\_batiments\_meme\_adresse[]} — jamais supprimées. Un test
        couvre le cas de trois bâtiments à une même adresse.
\end{enumerate}

\section{Les classes TRI : la profondeur d'eau réglementaire}

Pour la mise en situation sur la carte (chapitre~\ref{ch:scenarios}), il faut
une profondeur d'eau. La source utilisée est le \textbf{Territoire à Risque
d'Inondation} (TRI), issu de la Directive Inondation. Le TRI classe
l'exposition en quatre niveaux correspondant à des périodes de retour :

\begin{table}[H]
\centering
\small
\caption{Les classes TRI officielles et leur signification.}
\label{tab:tri}
\begin{tabular}{@{}l l l@{}}
\toprule
\textbf{Classe} & \textbf{Période de retour} & \textbf{Lecture} \\
\midrule
\code{frequent} & $\approx$ 10 ans & Événement courant \\
\code{moyen} & $\approx$ 100 ans & Événement de référence \\
\code{extreme} & $\approx$ 500 ans & Événement rare \\
\code{faible} & au-delà & Exposition résiduelle \\
\bottomrule
\end{tabular}
\end{table}

Le point crucial : ces classes et leurs profondeurs sont des \textbf{facts
réglementaires}. Elles ne sont jamais inventées ni interpolées. Si le point
n'est pas cartographié, le produit le dit — il ne devine pas une hauteur.

\section{Les sources complémentaires}

\begin{description}[leftmargin=0em,style=nextline]
  \item[Géocodage — Géoplateforme (IGN).] L'adresse saisie est convertie en
  coordonnées, code INSEE, code postal et libellé normalisé, avec un score de
  confiance. Le géocodage est traité comme un composant produit, pas comme une
  utilité : gestion des arrondissements, rejet des scores faibles, gestion de
  la saturation (HTTP 429 avec en-tête \code{retry-after}).
  \item[Hydrographie — Sandre.] Le tracé des cours d'eau, utilisé pour
  l'animation de suivi de l'eau sur la carte.
  \item[Météo — Open-Meteo.] Pluie horaire et débit GloFAS. Utilisée pour la
  \emph{forme} de la montée des eaux, jamais pour son niveau
  (voir chapitre~\ref{ch:scenarios}).
  \item[Photographie — Panoramax.] Une photo de terrain orientée sur le point,
  utile au souscripteur pour un contrôle visuel rapide.
  \item[Copernicus CDS.] Les projections climatiques à 2050 et 2100, en vraies
  unités (jours de canicule par an, degrés, mètres par seconde), pour les deux
  scénarios RCP 4.5 et RCP 8.5.
  \item[Mistral.] La rédaction du rapport narratif. Strictement confinée à la
  couche de présentation : elle ne produit aucun fait et n'entre jamais dans
  le contrat canonique.
\end{description}

\section{Contraintes juridiques et quotas}

\subsection{Licence Ouverte 2.0}

Géorisques et la BDNB sont distribués sous \textbf{Licence Ouverte 2.0}
(Etalab). La réutilisation commerciale est explicitement autorisée, sans
limitation de durée ni de territoire. La seule obligation est
l'\textbf{attribution} : nom de la source et date de dernière mise à jour.

Cette obligation coïncide exactement avec la conception de provenance du
produit : la conformité juridique devient une fonctionnalité, pas une charge.
Chaque bloc de provenance du contrat porte une chaîne d'attribution du type :

\begin{center}
\code{Source Géorisques (BRGM / MTE) — données à jour au 2026-07-01}
\end{center}

Le millésime est stocké en configuration par source et mis à jour lorsque le
fournisseur publie une nouvelle version.

\subsection{Quotas}

Les quotas des fournisseurs sont respectés côté client, car les dépasser
signifie une coupure de service :

\begin{table}[H]
\centering
\small
\caption{Quotas amont et mécanismes de protection implémentés.}
\label{tab:quotas}
\begin{tabularx}{\textwidth}{@{}l l X@{}}
\toprule
\textbf{Fournisseur} & \textbf{Quota} & \textbf{Protection} \\
\midrule
BDNB & 120 req/min/IP, 10\,000 req/mois (offre gratuite) &
  Limitation de débit côté client + compteur budgétaire mensuel \\
\addlinespace
Géoplateforme & 50 req/s/IP &
  Garde de rafale + gestion du 429 avec réattente unique \\
\bottomrule
\end{tabularx}
\end{table}

Un garde budgétaire mensuel protège l'allocation BDNB : lorsque le compteur est
épuisé, les soumissions de lots sont rejetées avec un HTTP 429
\code{budget\_epuise}, tandis que les requêtes unitaires interactives
continuent de fonctionner. Augmenter le débit réel est une action commerciale
(souscrire à l'offre supérieure), pas une décision de code.

\section{Le cache}

Un cache en mémoire à durée de vie limitée (TTL, 24~h par défaut,
configurable) évite de re-interroger les sources pour une même adresse. La clé
est construite sur le triplet de géocodage normalisé et sur la version de
schéma :

\begin{center}
\code{libellé normalisé} + \code{latitude/longitude arrondies à 5 décimales} +
\code{code INSEE} + \code{schema\_version}
\end{center}

L'arrondi à 5 décimales correspond à environ 1~m : deux adresses distinctes ne
partagent donc pas une clé par accident. Un point d'honnêteté : un résultat
servi depuis le cache conserve ses horodatages d'origine. Il indique la date de
la collecte réelle, pas la date de la consultation.

\begin{honesty}
Sur l'instance serverless déployée, ce cache ne persiste pas — le système de
fichiers de la fonction est éphémère. Chaque rapport narratif consomme donc un
appel au modèle, et la prose varie légèrement entre deux requêtes identiques.
En local, le cache rejoue le rapport à l'identique. Cette différence est
documentée dans le README plutôt que masquée.
\end{honesty}
